\documentclass[11pt]{article}
\usepackage[T1]{fontenc}
\usepackage[utf8]{inputenc}
\usepackage{lmodern,microtype}
\usepackage[margin=1in]{geometry}
\usepackage{amsmath,amssymb,amsthm,mathtools}
\usepackage{booktabs,array}
\usepackage[dvipsnames]{xcolor}
\usepackage{hyperref}
\usepackage{fancyhdr}
\hypersetup{colorlinks=true,linkcolor=MidnightBlue,citecolor=MidnightBlue,
 urlcolor=MidnightBlue,pdftitle={Three OEIS Conjectural Families for Apery Numbers},
 pdfauthor={OpenAI},pdfsubject={Finite telescoping, shifted diagonals, and zeta acceleration}}
\pagestyle{fancy}
\fancyhf{}
\fancyhead[L]{\small Shifted Ap\'ery series}
\fancyhead[R]{\small Finite telescoping and exact error bounds}
\fancyfoot[C]{\thepage}
\setlength{\headheight}{14pt}
\setlength{\emergencystretch}{2em}
\numberwithin{equation}{section}
\newtheorem{theorem}{Theorem}[section]
\newtheorem{lemma}[theorem]{Lemma}
\newtheorem{proposition}[theorem]{Proposition}
\newtheorem{corollary}[theorem]{Corollary}
\theoremstyle{definition}
\newtheorem{definition}[theorem]{Definition}
\theoremstyle{remark}
\newtheorem{remark}[theorem]{Remark}
\newcommand{\N}{\mathbb N}
\newcommand{\Z}{\mathbb Z}
\newcommand{\Q}{\mathbb Q}
\newcommand{\dd}{\,\mathrm d}
\newcommand{\ee}{\mathrm e}
\newcommand{\Apery}{Ap\'ery}
\newcommand{\HS}[2]{H_{#1}^{(#2)}}
\newcommand{\ES}[1]{E_{#1}^{(2)}}
\newcommand{\Sthree}[2]{S^{(3)}_{#1,#2}}
\newcommand{\Stwop}[2]{S^{(2,+)}_{#1,#2}}
\newcommand{\Stwom}[2]{S^{(2,-)}_{#1,#2}}
\DeclareMathOperator{\Li}{Li}

\title{\textbf{Three OEIS Conjectural Families\\for Ap\'ery Numbers}\\[5pt]
\large Shifted-diagonal zeta series from finite telescoping}
\author{A self-contained mathematical investigation}
\date{September 20, 2026}

\begin{document}
\maketitle
\begin{abstract}
We prove three entire families of series identities explicitly labeled conjectural
in OEIS entries A143007 and A108625 when consulted on September 20, 2026.
The first expresses $\zeta(3)$ through every shifted diagonal of a symmetric
binomial array whose main diagonal is the classical \Apery{} sequence
$1,5,73,1445,\ldots$. The other two express $\zeta(2)$ through the superdiagonals
and subdiagonals of the array with main diagonal $1,3,19,147,\ldots$.
All nonnegative integer shifts are covered, not merely individual numerical
instances. Explicit finite telescoping certificates establish local relations
between four neighboring array entries. These relations produce path-independent
rational sums on the integer lattice, from which the conjectured formulas and
exact remainder identities follow. We also derive certified rational error
bounds, shifted recurrences, and sharp leading asymptotics for all three
families. Supplementary programs perform 25,840 exact-arithmetic checks and
construct rational enclosures certifying at least 120 decimal places.
The proofs are independent of numerical evidence. No claim of historical
priority is made; the relationship with existing conservative-matrix-field
work is identified explicitly.
\end{abstract}

\newpage
\tableofcontents
\newpage

\section{The conjectures and their mathematical setting}\label{sec:main}

\subsection{Why these sequences?}
The two classical \Apery{} sequences are
\begin{align}
 b_n&=\sum_{j=0}^n\binom nj^2\binom{n+j}j
       =1,3,19,147,1251,\ldots,\label{eq:b}\\
 a_n&=\sum_{j=0}^n\binom nj^2\binom{n+j}j^2
       =1,5,73,1445,33001,\ldots.\label{eq:a}
\end{align}
They are OEIS A005258 and A005259, respectively~\cite{oeis258,oeis259}.
Their connection with rational approximation to $\zeta(2)$ and $\zeta(3)$
is central to \Apery{}'s work~\cite{apery,david}. The problems considered here
concern two-dimensional extensions of these sequences, rather than an
isolated numerical pattern in an unrelated sequence.

Define, for $n,m\in\Z_{\geq0}$,
\begin{align}
 U(n,m)&=\sum_{j=0}^{\min(n,m)}
      \binom nj\binom{n+j}j\binom mj,\label{eq:U}\\
 T(n,m)&=\sum_{j=0}^{\min(n,m)}
      \binom nj\binom{n+j}j\binom mj\binom{m+j}j.\label{eq:T}
\end{align}
These are the square arrays A108625 and A143007~\cite{oeisU,oeisT}.
In particular, $U(n,n)=b_n$ and $T(n,n)=a_n$.
The OEIS stores each array as a single sequence read by antidiagonals;
throughout this article, the notation denotes the square-array coordinates,
not the flattened indexing. Both OEIS entries use the letter $T$ for their
own array; we use $U$ for A108625 to avoid confusion.

Put
\[
 \HS{k}{s}=\sum_{r=1}^k\frac1{r^s},\qquad
 \ES{k}=\sum_{r=1}^k\frac{(-1)^{r+1}}{r^2},
 \qquad \HS{0}{s}=\ES{0}=0.
\]
Only $s=2,3$ will be needed, and $\zeta(s)=\sum_{r\geq1}r^{-s}$.

\subsection{Three full families}
The following statements are the principal results.

\begin{theorem}[Shifted diagonals for $\zeta(3)$]\label{thm:main3}
For every integer $k\geq0$,
\begin{equation}\label{eq:main3}
\boxed{\displaystyle
\zeta(3)=\HS{k}{3}
 +\sum_{n=1}^{\infty}
 \frac{(2n+k)(3n^2+3nk+k^2)}
 {n^3(n+k)^3T(n-1,n+k-1)T(n,n+k)}.}
\end{equation}
\end{theorem}

\begin{theorem}[Superdiagonals for $\zeta(2)$]\label{thm:main2p}
For every integer $k\geq0$,
\begin{equation}\label{eq:main2p}
\boxed{\displaystyle
\zeta(2)=\HS{k}{2}
 +\sum_{n=1}^{\infty}(-1)^{n+1}
 \frac{5n^2+6kn+2k^2}
 {n^2(n+k)^2U(n-1,n+k-1)U(n,n+k)}.}
\end{equation}
\end{theorem}

\begin{theorem}[Subdiagonals for $\zeta(2)$]\label{thm:main2m}
For every integer $k\geq0$,
\begin{equation}\label{eq:main2m}
\boxed{\displaystyle
\zeta(2)=2\ES{k}
 +(-1)^k\sum_{n=1}^{\infty}(-1)^{n+1}
 \frac{5n^2+4kn+k^2}
 {n^2(n+k)^2U(n+k-1,n-1)U(n+k,n)}.}
\end{equation}
\end{theorem}

Every series above is absolutely convergent. The $k=0$ cases are the familiar
main-diagonal formulas. The assertions for \emph{arbitrary shifts} are the
conjectural families targeted here.

\subsection{Source status and scope of the claim}
At the date of consultation, the formula field of A143007 labels
\eqref{eq:main3} a ``Conjectural result for other diagonals''~\cite{oeisT}.
The formula field of A108625 separately labels the two families
\eqref{eq:main2p} and \eqref{eq:main2m} conjectural for superdiagonals and
subdiagonals~\cite{oeisU}. The article proves exactly those formulas, with
all signs, initial harmonic sums, and coordinate shifts retained.

An OEIS label is evidence of the status of that entry, not proof that a
formula has never appeared elsewhere. In particular, David's
conservative-matrix-field paper already develops related lattice
constructions for these constants and contains, up to its duality and sign conventions, the conjugate
polynomials occurring below~\cite[Examples 16, 19, 25; Section 5]{david}.
Thus our claim is a self-contained proof of formulas \emph{listed as
conjectural}, not a claim to have originated the underlying framework.
The arithmetic theory of multivariate \Apery{} arrays is also developed
in Straub's work~\cite{straub}. No supercongruence or irreducibility
conjecture is asserted to be settled here.

\section{Elementary facts and a lattice-sum lemma}\label{sec:lattice}

Each term in \eqref{eq:U} and \eqref{eq:T} is a nonnegative integer, and the
$j=0$ term is $1$. Thus both arrays are positive and their entries can safely
appear in denominators. Directly from the sums,
\begin{gather}
 U(n,0)=U(0,m)=T(n,0)=T(0,m)=1,\label{eq:axes}\\
 T(n,m)=T(m,n),\label{eq:symmetry}\\
 T(n,1)=1+2n(n+1),\quad
 U(n,1)=1+n(n+1),\quad U(1,m)=1+2m.\label{eq:firstrow}
\end{gather}
Both arrays are nondecreasing in either coordinate: when a coordinate
increases, every existing binomial factor is nondecreasing, and any
new terms are nonnegative. Only $T$ is symmetric; for example,
$U(1,2)=5$ whereas $U(2,1)=7$.

\begin{lemma}[Closed edge weights give a potential]\label{lem:closed}
Assign a rational weight $h(n,m)$ to the east edge
$(n-1,m)\to(n,m)$, for $n\geq1,m\geq0$, and a rational weight $v(n,m)$
to the north edge $(n,m-1)\to(n,m)$, for $n\geq0,m\geq1$.
Suppose that for every $n,m\geq1$,
\begin{equation}\label{eq:closedgeneral}
 h(n,m-1)+v(n,m)=v(n-1,m)+h(n,m).
\end{equation}
Then the sum along a monotone east/north path from $(0,0)$ to $(n,m)$
depends only on the endpoint. Denote it by $F(n,m)$, with $F(0,0)=0$.
Its horizontal and vertical differences are precisely $h$ and $v$.
\end{lemma}
\begin{proof}
A path is a word with $n$ letters E and $m$ letters N. Any two such words
are related by finitely many interchanges EN$\leftrightarrow$NE of adjacent
unlike letters: both can, for instance, be sorted to E$^n$N$^m$.
Each interchange replaces two sides of a unit square by its other two
sides, and \eqref{eq:closedgeneral} preserves the sum. Appending a single
edge proves the stated difference formulas.
\end{proof}

This lemma is elementary. The real work is to find and prove the
four-entry identities which make the chosen weights closed, and then
to identify the limiting value of the resulting potential. A local
identity alone does not determine an infinite sum.

\section{A finite certificate for the cubic array}\label{sec:cubic}

Fix $n,m\geq1$ and name the corners of one square as follows:
\[
 A=T(n-1,m-1),\quad B=T(n-1,m),\quad
 C=T(n,m-1),\quad D=T(n,m).
\]
Thus $A,B,C,D$ are respectively the lower-left, upper-left, lower-right,
and upper-right entries. Define
\begin{align}
 f_3(n,m)&=(n+m)(n^2+nm+m^2),\label{eq:f3}\\
 g_3(n,m)&=(n-m)(n^2-nm+m^2).\label{eq:g3}
\end{align}
In particular,
\begin{equation}\label{eq:fg3}
 f_3(n,m)g_3(n,m)=n^6-m^6.
\end{equation}

\begin{proposition}[Cubic contiguity relations]\label{prop:contig3}
The four corners satisfy
\begin{equation}\label{eq:contig3}
 \boxed{f_3B=n^3D+m^3A,\qquad f_3C=m^3D+n^3A.}
\end{equation}
These identities remain valid when $n=m$; no division by $n-m$ is used.
\end{proposition}

\begin{proof}
Let $q=\min(n,m)$ and
\[
 c_j=\binom nj\binom{n+j}j\binom mj\binom{m+j}j
 \quad(0\leq j\leq q).
\]
Relative to the summand $c_j$ of $D$, the summands of $B,C,A$ have ratios
\begin{equation}\label{eq:corner-ratios3}
 b_j=\frac{n-j}{n+j},\qquad
 c'_j=\frac{m-j}{m+j},\qquad
 a_j=b_jc'_j.
\end{equation}
These formulas include a vanished boundary term when $j=n$ or $j=m$.
The rational identity
\[
 (n+m)(c'_j-b_j)=(n-m)(a_j-1)
\]
gives, after multiplication by $c_j$ and summation,
\begin{equation}\label{eq:linear3}
 (n+m)(C-B)=(n-m)(A-D).
\end{equation}

For the second relation, define the explicit certificate
\begin{equation}\label{eq:W}
 W_j=\frac{j^4c_j}{(n+j)(m+j)}\quad(0\leq j\leq q),
 \qquad W_{q+1}=0.
\end{equation}
We have $W_0=0$. The ratio of consecutive summands is
\[
 \frac{c_{j+1}}{c_j}
 =\frac{(n-j)(n+j+1)(m-j)(m+j+1)}{(j+1)^4}.
\]
Consequently, for $0\leq j\leq q$,
\begin{equation}\label{eq:deltaW}
 W_{j+1}-W_j
 =\frac{n^2m^2-(n^2+m^2)j^2}{(n+j)(m+j)}c_j.
\end{equation}
At $j=q$ the same equation holds because one of $n-j,m-j$ vanishes;
there is no unaccounted terminal term. Hence the sum of the right side
of \eqref{eq:deltaW} is zero.

Writing $h=n^2-nm+m^2$ and $h_+=n^2+nm+m^2$, direct expansion gives
\[
 \bigl(h(a_j+1)-h_+(b_j+c'_j)\bigr)c_j
 =-4(W_{j+1}-W_j).
\]
Summation proves
\begin{equation}\label{eq:quadratic3}
 h(A+D)=h_+(B+C).
\end{equation}
Equations \eqref{eq:linear3} and \eqref{eq:quadratic3} determine $B$ and $C$.
For clarity, multiplying the expression
\[
 2B=\frac{h}{h_+}(A+D)-\frac{n-m}{n+m}(A-D)
\]
by $f_3=(n+m)h_+$ yields
\[
 2f_3B=(n^3+m^3)(A+D)-(n^3-m^3)(A-D)
       =2m^3A+2n^3D.
\]
The formula for $C$ follows by changing the minus sign to a plus sign.
All denominators used are positive for $n,m\geq1$.
\end{proof}

\begin{corollary}[Cubic square closure]\label{cor:closed3}
The corners satisfy
\begin{equation}\label{eq:closed3}
 \frac1{n^3AC}+\frac1{m^3CD}
 =\frac1{m^3AB}+\frac1{n^3BD}.
\end{equation}
\end{corollary}
\begin{proof}
The relations \eqref{eq:contig3}, together with \eqref{eq:fg3}, give
\[
 n^3C-m^3B=g_3A,\qquad n^3B-m^3C=g_3D.
\]
Indeed, multiplying either claimed relation by $f_3$ and using
\eqref{eq:contig3} reduces it to \eqref{eq:fg3}.
Therefore
$D(n^3C-m^3B)=A(n^3B-m^3C)$.
Rearranging and dividing by the positive number $n^3m^3ABCD$
gives \eqref{eq:closed3}.
\end{proof}

\section{The cubic potential and the first conjecture}\label{sec:proof3}

Define
\begin{align}
 h_3(n,m)&=\frac1{n^3T(n-1,m)T(n,m)},\label{eq:h3}\\
 v_3(n,m)&=\frac1{m^3T(n,m-1)T(n,m)}.\label{eq:v3}
\end{align}
Corollary~\ref{cor:closed3} is exactly the hypothesis of
Lemma~\ref{lem:closed}. Let $F_3(n,m)$ be the resulting potential.
Taking first $n$ east steps on the axis and then $m$ north steps gives
\begin{equation}\label{eq:F3}
 F_3(n,m)=\HS{n}{3}+
 \sum_{j=1}^m\frac1{j^3T(n,j-1)T(n,j)}.
\end{equation}
Taking the transposed path and using symmetry of $T$ shows that
\begin{equation}\label{eq:F3sym}
 F_3(n,m)=F_3(m,n).
\end{equation}
Every edge weight is positive, so $F_3$ increases strictly when either
coordinate is increased.

\begin{proposition}[The limiting value]\label{prop:limit3}
For every finite $n,m$,
\begin{equation}\label{eq:F3sandwich}
 \HS{\max(n,m)}{3}\leq F_3(n,m)<\zeta(3).
\end{equation}
Moreover, $F_3(n,m)\to\zeta(3)$ whenever $\max(n,m)\to\infty$.
\end{proposition}
\begin{proof}
Fix $m$. In each of the finitely many summands in \eqref{eq:F3},
\[
 T(n,j-1)T(n,j)\geq T(n,1)=1+2n(n+1)\longrightarrow\infty.
\]
Thus $F_3(n,m)\to\zeta(3)$ as $n\to\infty$ with $m$ fixed.
Since $F_3$ is strictly increasing in $n$, every finite value is strictly
below that limit. Formula \eqref{eq:F3} and its transposed version give
the lower bound in \eqref{eq:F3sandwich}. Letting the larger coordinate
tend to infinity now proves the last assertion by squeezing.
\end{proof}

\begin{proposition}[One diagonal step]\label{prop:step3}
For $n,m\geq1$,
\begin{equation}\label{eq:step3}
 F_3(n,m)-F_3(n-1,m-1)
 =\frac{f_3(n,m)}
 {n^3m^3T(n-1,m-1)T(n,m)}.
\end{equation}
\end{proposition}
\begin{proof}
Use the east-then-north path across the square. In the corner notation,
the increment is
\[
 \frac1{n^3AC}+\frac1{m^3CD}
 =\frac{m^3D+n^3A}{n^3m^3ACD}.
\]
The second relation in \eqref{eq:contig3} replaces the numerator by $f_3C$,
which cancels the intermediate corner.
\end{proof}

\begin{proof}[Proof of Theorem~\ref{thm:main3}]
Put $(n,m)=(r,r+k)$ in \eqref{eq:step3} and sum from $r=1$ to $r=N$.
The left side telescopes to $F_3(N,N+k)-F_3(0,k)$.
Since $F_3(0,k)=\HS{k}{3}$ and
\[
 f_3(r,r+k)=(2r+k)(3r^2+3rk+k^2),
\]
the finite identity is exactly
\begin{equation}\label{eq:finite3}
 \Sthree{N}{k}:=\HS{k}{3}+
 \sum_{r=1}^N\frac{(2r+k)(3r^2+3rk+k^2)}
 {r^3(r+k)^3T(r-1,r+k-1)T(r,r+k)}=F_3(N,N+k).
\end{equation}
Proposition~\ref{prop:limit3} allows $N\to\infty$ and gives
\eqref{eq:main3}. Positivity makes this a convergent positive series.
\end{proof}

\subsection{Exact cubic remainder identities}
\begin{theorem}[Remainder and certified bounds]\label{thm:error3}
For all $n,m\geq0$,
\begin{align}
 \zeta(3)-F_3(n,m)
 &=\sum_{r=n+1}^\infty
   \frac1{r^3T(r-1,m)T(r,m)}\label{eq:tail3h}\\
 &=\sum_{j=m+1}^\infty
   \frac1{j^3T(n,j-1)T(n,j)}.\label{eq:tail3v}
\end{align}
If $M=\max(n,m)\geq1$, then
\begin{equation}\label{eq:bound3}
 0<\zeta(3)-F_3(n,m)<\frac1{2M^2T(n,m)^2}.
\end{equation}
In particular, for $N\geq1$ and $k\geq0$,
\begin{equation}\label{eq:bound3diag}
 0<\zeta(3)-\Sthree{N}{k}
 <\frac1{2(N+k)^2T(N,N+k)^2}.
\end{equation}
\end{theorem}
\begin{proof}
Telescope horizontal edge weights from $(n,m)$ to $(L,m)$ and let
$L\to\infty$, using Proposition~\ref{prop:limit3}. This proves
\eqref{eq:tail3h}; the vertical version is identical.
Use the tail whose lower index is the larger coordinate. Every array
factor in that tail is at least $T(n,m)$ by monotonicity, and
\[
 \sum_{r=M+1}^\infty r^{-3}<\int_M^\infty x^{-3}\dd x=\frac1{2M^2}.
\]
This proves \eqref{eq:bound3}, and \eqref{eq:finite3} gives the diagonal case.
\end{proof}

The first omitted positive diagonal summand is also a strict lower
bound for the error. Thus an interval can be narrowed on both sides
using only rational arithmetic. The bound is a theorem, not an estimate
obtained by comparing with a stored numerical value of $\zeta(3)$.

\section{A finite certificate for the quadratic array}\label{sec:quadratic}

Now use the analogous corner notation for $U$:
\[
 A=U(n-1,m-1),\quad B=U(n-1,m),\quad
 C=U(n,m-1),\quad D=U(n,m),
 \qquad n,m\geq1.
\]
Set
\begin{equation}\label{eq:f2g2}
 f_2=n^2+2nm+2m^2,\qquad g_2=n^2-2nm+2m^2.
\end{equation}
These satisfy
\begin{equation}\label{eq:fg2}
 f_2g_2=n^4+4m^4.
\end{equation}

\begin{proposition}[Quadratic contiguity relations]\label{prop:contig2}
For every $n,m\geq1$,
\begin{equation}\label{eq:contig2}
 \boxed{f_2B=n^2D+2m^2A,\qquad f_2C=2m^2D-n^2A.}
\end{equation}
\end{proposition}
\begin{proof}
Let $q=\min(n,m)$ and
\[
 d_j=\binom nj\binom{n+j}j\binom mj\quad(0\leq j\leq q).
\]
The ratios of the summands of $B,C,A$ to $d_j$ are now
\[
 b_j=\frac{n-j}{n+j},\qquad c'_j=\frac{m-j}{m},\qquad a_j=b_jc'_j.
\]
Define
\begin{equation}\label{eq:V}
 V_j=\frac{j^3d_j}{n+j}\quad(0\leq j\leq q),\qquad V_{q+1}=0.
\end{equation}
As before $V_0=0$, and the summand ratio
\[
 \frac{d_{j+1}}{d_j}
 =\frac{(n-j)(n+j+1)(m-j)}{(j+1)^3}
\]
gives the identity, including the endpoint $j=q$,
\begin{equation}\label{eq:deltaV}
 V_{j+1}-V_j=\frac{mn^2-n^2j-mj^2}{n+j}d_j.
\end{equation}
Direct expansion yields the two termwise equalities
\begin{align*}
 (f_2b_j-n^2-2m^2a_j)d_j&=2(V_{j+1}-V_j),\\
 (f_2c'_j-2m^2+n^2a_j)d_j&=\frac{2(n+m)}m(V_{j+1}-V_j).
\end{align*}
Summing and using both zero boundary values proves \eqref{eq:contig2}.
\end{proof}

\begin{corollary}[Quadratic square closure]\label{cor:closed2}
The corners satisfy
\begin{equation}\label{eq:closed2}
 \frac2{n^2AC}-\frac1{m^2CD}
 =\frac1{m^2AB}+\frac2{n^2BD}.
\end{equation}
\end{corollary}
\begin{proof}
Equations \eqref{eq:contig2} and \eqref{eq:fg2} imply
\[
 2m^2B-n^2C=g_2A,\qquad n^2B+2m^2C=g_2D.
\]
Multiplying the first by $D$, the second by $A$, and equating gives
\[
 D(2m^2B-n^2C)=A(n^2B+2m^2C).
\]
Divide by $n^2m^2ABCD$ to obtain the result.
\end{proof}

The minus sign in \eqref{eq:closed2} is essential. Unlike the cubic case,
this construction requires signed edge weights; treating $U$ as a
symmetric array or dropping an alternating sign gives a false formula.

\section{The signed potential and both quadratic conjectures}\label{sec:proof2}

Assign the weights
\begin{align}
 h_2(n,m)&=\frac{2(-1)^{n+1}}{n^2U(n-1,m)U(n,m)},\label{eq:h2}\\
 v_2(n,m)&=\frac{(-1)^n}{m^2U(n,m-1)U(n,m)}.\label{eq:v2}
\end{align}
After division by the common sign $(-1)^{n+1}$, square closure becomes
exactly \eqref{eq:closed2}. Lemma~\ref{lem:closed} therefore gives a
rational potential $F_2(n,m)$. Its axis-then-vertical expression is
\begin{equation}\label{eq:F2}
 F_2(n,m)=2\ES{n}
 +(-1)^n\sum_{j=1}^m\frac1{j^2U(n,j-1)U(n,j)}.
\end{equation}
Its boundary values are
\begin{equation}\label{eq:F2axes}
 F_2(n,0)=2\ES{n},\qquad F_2(0,m)=\HS{m}{2}.
\end{equation}

\begin{theorem}[Quadratic limit and exact errors]\label{thm:error2}
The potential tends to $\zeta(2)$ whenever $\max(n,m)\to\infty$.
For all finite $n,m\geq0$,
\begin{align}
 \zeta(2)-F_2(n,m)
 &=2\sum_{r=n+1}^{\infty}
 \frac{(-1)^{r+1}}{r^2U(r-1,m)U(r,m)}\label{eq:tail2h}\\
 &=(-1)^n\sum_{j=m+1}^{\infty}
 \frac1{j^2U(n,j-1)U(n,j)}.\label{eq:tail2v}
\end{align}
In particular,
\begin{equation}\label{eq:bound2h}
 0<(-1)^n\bigl(\zeta(2)-F_2(n,m)\bigr)
 <\frac2{(n+1)^2U(n,m)U(n+1,m)}.
\end{equation}
For $m\geq1$, the additional bound
\begin{equation}\label{eq:bound2v}
 0<(-1)^n\bigl(\zeta(2)-F_2(n,m)\bigr)
 <\frac1{mU(n,m)^2}
\end{equation}
holds. The smaller of the two rational upper bounds may be used.
\end{theorem}

\begin{proof}
First fix $m$. Since $U(n,j-1)U(n,j)\geq U(n,1)\to\infty$,
the finite correction in \eqref{eq:F2} tends to zero as $n\to\infty$.
Absolute convergence permits the elementary even/odd decomposition
\[
 \sum_{r=1}^\infty\frac{(-1)^{r+1}}{r^2}
 =\sum_{r=1}^\infty\frac1{r^2}
   -2\sum_{r=1}^\infty\frac1{(2r)^2}
 =\frac12\zeta(2).
\]
Thus $F_2(n,m)\to\zeta(2)$ for each fixed $m$.

Telescoping the horizontal edges and taking this limit proves
\eqref{eq:tail2h}. The positive numbers
\[
 \frac{2}{r^2U(r-1,m)U(r,m)}
\]
strictly decrease to zero as $r$ increases: the factor $r^2$ increases
strictly, and the array factors are positive and nondecreasing.
The alternating-series remainder lies strictly between zero and its
first term, with the sign of that term. At $r=n+1$ the sign is $(-1)^n$.
This proves \eqref{eq:bound2h}, including its strict sign assertion.

The upper bound also proves convergence for arbitrary diverging
endpoints, not just for a fixed row. Indeed, monotonicity gives
$U(n+1,m)\geq U(1,m)=2m+1$ and $U(n,m)\geq1$, so
\begin{equation}\label{eq:uniform2}
 |\zeta(2)-F_2(n,m)|
 <\frac2{(n+1)^2(2m+1)}.
\end{equation}
The denominator tends to infinity whenever $\max(n,m)$ does.
In particular, with $n$ fixed we may let $m\to\infty$ in a telescoped
vertical path, obtaining \eqref{eq:tail2v}. Its terms have a common sign.
For $m\geq1$, monotonicity and
$\sum_{j>m}j^{-2}<\int_m^\infty x^{-2}\dd x=1/m$
give \eqref{eq:bound2v}.
\end{proof}

\begin{proposition}[One quadratic diagonal step]\label{prop:step2}
For $n,m\geq1$,
\begin{equation}\label{eq:step2}
 F_2(n,m)-F_2(n-1,m-1)
 =(-1)^{n+1}\frac{n^2+2nm+2m^2}
 {n^2m^2U(n-1,m-1)U(n,m)}.
\end{equation}
\end{proposition}
\begin{proof}
Following an east edge and then a north edge across the square gives
\[
 (-1)^{n+1}\left(\frac2{n^2AC}-\frac1{m^2CD}\right)
 =(-1)^{n+1}\frac{2m^2D-n^2A}{n^2m^2ACD}.
\]
Use $2m^2D-n^2A=f_2C$ from \eqref{eq:contig2} and cancel $C$.
\end{proof}

\begin{proof}[Proof of Theorems~\ref{thm:main2p} and \ref{thm:main2m}]
For the superdiagonal set $(n,m)=(r,r+k)$ in \eqref{eq:step2}.
The numerator becomes $5r^2+6kr+2k^2$, the sign is $(-1)^{r+1}$,
and the initial point $(0,k)$ contributes $\HS{k}{2}$.
For the subdiagonal set $(n,m)=(r+k,r)$. The numerator is
$5r^2+4kr+k^2$, the sign is $(-1)^{r+k+1}$, and the initial point
$(k,0)$ contributes $2\ES{k}$.
Thus the finite partial sums of \eqref{eq:main2p} and
\eqref{eq:main2m} satisfy, respectively,
\begin{equation}\label{eq:finite2}
 \Stwop{N}{k}=F_2(N,N+k),\qquad
 \Stwom{N}{k}=F_2(N+k,N).
\end{equation}
The limiting assertion in Theorem~\ref{thm:error2} proves both infinite
identities. Absolute convergence follows from the elementary lower
bounds established in Section~\ref{sec:examples} below.
\end{proof}

\begin{corollary}[Certified quadratic diagonal errors]\label{cor:diag2}
For $N\geq1,k\geq0$,
\begin{align}
0&<(-1)^N\bigl(\zeta(2)-\Stwop{N}{k}\bigr)\notag\\
 &<\min\left\{
 \frac2{(N+1)^2U(N,N+k)U(N+1,N+k)},
 \frac1{(N+k)U(N,N+k)^2}\right\},\label{eq:diagerror2p}\\
0&<(-1)^{N+k}\bigl(\zeta(2)-\Stwom{N}{k}\bigr)\notag\\
 &<\min\left\{
 \frac2{(N+k+1)^2U(N+k,N)U(N+k+1,N)},
 \frac1{NU(N+k,N)^2}\right\}.\label{eq:diagerror2m}
\end{align}
\end{corollary}
\begin{proof}
Substitute the two endpoints in \eqref{eq:bound2h} and \eqref{eq:bound2v}.
\end{proof}

\section{Examples, absolute convergence, and nontrivial acceleration}\label{sec:examples}

At $k=0$, the theorems reduce to
\begin{equation}\label{eq:classicalseries}
 \zeta(3)=6\sum_{n=1}^\infty\frac1{n^3a_{n-1}a_n},\qquad
 \zeta(2)=5\sum_{n=1}^\infty\frac{(-1)^{n+1}}{n^2b_{n-1}b_n}.
\end{equation}
The initial partial sums give a useful check on normalization:
\begin{center}
\begin{tabular}{@{}rcc@{}}
\toprule
$N$ & $\Sthree{N}{0}$ & $\Stwop{N}{0}=\Stwom{N}{0}$\\
\midrule
$1$ & $6/5$ & $5/3$\\
$2$ & $351/292$ & $125/76$\\
$3$ & $62531/52020$ & $8705/5292$\\
\bottomrule
\end{tabular}
\end{center}

At shift $k=1$, the cubic identity is
\[
 \zeta(3)=1+\sum_{n\geq1}
 \frac{(2n+1)(3n^2+3n+1)}
 {n^3(n+1)^3T(n-1,n)T(n,n+1)}.
\]
The quadratic pair is
\begin{align*}
 \zeta(2)&=1+\sum_{n\geq1}(-1)^{n+1}
 \frac{5n^2+6n+2}{n^2(n+1)^2U(n-1,n)U(n,n+1)},\\
 \zeta(2)&=2+\sum_{n\geq1}(-1)^n
 \frac{5n^2+4n+1}{n^2(n+1)^2U(n,n-1)U(n+1,n)}.
\end{align*}
Their first approximants are respectively $125/104$, $33/20$, and $23/14$.
The different signs and different boundary terms in the quadratic pair
are consequences of the asymmetric signed potential.

\begin{proposition}[Elementary exponential bounds]\label{prop:coarse}
For every $N\geq1,k\geq0$,
\begin{align}
 T(N,N+k)&\geq\binom{2N}N^2\geq\frac{16^N}{(2N+1)^2},\label{eq:coarseT}\\
 U(N,N+k),\ U(N+k,N)&\geq\binom{2N}N\geq\frac{4^N}{2N+1}.\label{eq:coarseU}
\end{align}
Consequently,
\begin{align}
0<\zeta(3)-\Sthree{N}{k}
 &<\frac{(2N+1)^4}{2(N+k)^2\,256^N},\label{eq:coarseerrT}\\
|\zeta(2)-\Stwop{N}{k}|&<\frac{(2N+1)^2}{(N+k)16^N},\label{eq:coarseerrUp}\\
|\zeta(2)-\Stwom{N}{k}|&<\frac{(2N+1)^2}{N16^N}.\label{eq:coarseerrUm}
\end{align}
\end{proposition}
\begin{proof}
By coordinate monotonicity it suffices to use the main diagonal.
In the defining sums, the $j=N$ term is $\binom{2N}N^2$ for $T$ and
$\binom{2N}N$ for $U$. The largest of the $2N+1$ binomial coefficients
$\binom{2N}j$ is the central one, whereas their sum is $4^N$.
This proves the array bounds. Substitution in
\eqref{eq:bound3diag}, \eqref{eq:diagerror2p}, and
\eqref{eq:diagerror2m} gives the remainder bounds.
\end{proof}

For fixed $k$, the numerator divided by the explicit powers of $n,n+k$
in a cubic summand is $O_k(n^{-3})$, and in a quadratic summand it is
$O_k(n^{-2})$. Combining this with \eqref{eq:coarseT} or
\eqref{eq:coarseU} at $n$ and $n-1$ proves absolute convergence of all
three series. These bounds already establish exponential acceleration,
without requiring any asymptotic analysis. In contrast, the defining
partial sums of $\zeta(3)$ and $\zeta(2)$ have errors of orders $N^{-2}$
and $N^{-1}$, respectively, by the integral comparison for decreasing
powers.

\section{Recurrences and rational companion solutions}\label{sec:recurrences}

The finite identities also explain why \Apery{}-type three-term
recurrences accompany these sums. We give the cubic construction in
detail, including every fixed shift.

\begin{proposition}[A row recurrence]\label{prop:rowrec}
For $n\geq1,m\geq0$,
\begin{equation}\label{eq:rowrec}
 (n+1)^3T(n+1,m)
 =(2n+1)(n^2+n+2m^2+2m+1)T(n,m)-n^3T(n-1,m).
\end{equation}
\end{proposition}
\begin{proof}
For $m\geq1$, apply the first relation in \eqref{eq:contig3} to the
square with upper-right corner $(n+1,m)$:
\[
 f_3(n+1,m)T(n,m)=(n+1)^3T(n+1,m)+m^3T(n,m-1).
\]
The conjugate relation from the preceding square gives
\[
 m^3T(n,m-1)=n^3T(n-1,m)-g_3(n,m)T(n,m).
\]
Eliminate $T(n,m-1)$ and expand
\[
 f_3(n+1,m)+g_3(n,m)
 =(2n+1)(n^2+n+2m^2+2m+1).
\]
For $m=0$, every array entry is $1$ and the recurrence reduces to the
identity $(n+1)^3=(2n+1)(n^2+n+1)-n^3$.
\end{proof}

Fix $k\geq0$, and define
\begin{gather}
 q_n=T(n,n+k),\qquad f_n=f_3(n,n+k),\qquad p_n=n^3(n+k)^3,\label{eq:recdefs}\\
 L_n=(2n+1)\bigl(n^2+n+2(n+k)^2+2(n+k)+1\bigr),\nonumber\\
 K_n=f_{n+1}(f_nL_n-n^6)-f_n(n+1)^6.\label{eq:Kn}
\end{gather}
Here $f_n,p_n$ will be used only for $n\geq1$, so they are positive.

\begin{theorem}[A recurrence on each shifted diagonal]\label{thm:shiftrec}
For $n\geq1$,
\begin{equation}\label{eq:shiftrec}
 p_{n+1}f_nq_{n+1}=K_nq_n-f_{n+1}p_nq_{n-1},
\end{equation}
with
\begin{equation}\label{eq:qinit}
 q_0=1,\qquad q_1=2k^2+6k+5.
\end{equation}
The rational sequence
\begin{equation}\label{eq:rdef}
 r_n=q_n\Sthree{n}{k}
\end{equation}
satisfies the same recurrence, with
\begin{equation}\label{eq:rinit}
 r_0=\HS{k}{3},\qquad
 r_1=q_1\HS{k}{3}+\frac{f_1}{(k+1)^3}.
\end{equation}
\end{theorem}
\begin{proof}
Write $m=n+k$. The first contiguity relation expresses
\[
 T(n-1,m)=\frac{n^3q_n+m^3q_{n-1}}{f_n}.
\]
Substitute this in the row recurrence to get
\[
 (n+1)^3f_nT(n+1,m)
 =(f_nL_n-n^6)q_n-p_nq_{n-1}.
\]
The second contiguity relation at $(n+1,m+1)$ is
\[
 f_{n+1}T(n+1,m)=(m+1)^3q_{n+1}+(n+1)^3q_n.
\]
Eliminating $T(n+1,m)$ proves \eqref{eq:shiftrec}.
The initial values of $q$ follow from \eqref{eq:firstrow}.

Equation \eqref{eq:finite3} gives the discrete Wronskian
\begin{equation}\label{eq:wronskian}
 r_nq_{n-1}-r_{n-1}q_n=\frac{f_n}{p_n}\quad(n\geq1).
\end{equation}
Divide \eqref{eq:shiftrec} by $p_{n+1}f_n$ and write it as
$q_{n+1}=\alpha_nq_n-\beta_nq_{n-1}$, where
\[
 \beta_n=\frac{f_{n+1}p_n}{f_np_{n+1}}.
\]
If $W_n$ denotes the left side of \eqref{eq:wronskian}, then
$W_{n+1}=\beta_nW_n$. Therefore
\begin{align*}
 q_n(r_{n+1}-\alpha_nr_n+\beta_nr_{n-1})
 &=W_{n+1}-\beta_nW_n=0.
\end{align*}
Since $q_n>0$, the companion recurrence follows. Its initial values
come from the first diagonal increment.
\end{proof}

At $k=0$, simplification recovers the classical recurrence
\begin{equation}\label{eq:aperyrec}
 (n+1)^3a_{n+1}
 =(2n+1)(17n^2+17n+5)a_n-n^3a_{n-1},
\end{equation}
and the rational companion starts with $r_0=0,r_1=6$.
For computation along a fixed shift, \eqref{eq:shiftrec} avoids evaluating
an entire two-dimensional array. It requires $O(N)$ recurrence steps
to reach index $N$; this is an arithmetic-operation count, not a
unit-cost claim for arbitrarily large integers.

There is an important arithmetic distinction: $q_n$ is an integer,
but $r_n$ need not be an integer. Thus the approximation estimate alone
is not, without an additional denominator argument, a self-contained
new proof of irrationality. We do not conflate fast approximation with
that stronger conclusion.

\section{Sharp asymptotics for all three families}\label{sec:asymptotics}

We now go beyond bounds and calculate the exact leading error constants.
All asymptotic statements in this section have \emph{fixed} $k$ and
$n$ or $N$ tending to infinity. No uniform estimate for growing $k$
is asserted.

Put
\begin{equation}\label{eq:constants}
 \beta=3+2\sqrt2,\qquad \lambda=\beta^2=17+12\sqrt2,
 \qquad \varphi=\frac{1+\sqrt5}{2},\qquad
 \rho=\varphi^5=\frac{11+5\sqrt5}{2}.
\end{equation}

\begin{theorem}[Asymptotics of the array entries]\label{thm:arrayasy}
For each fixed $k\geq0$,
\begin{align}
 T(n,n+k)&\sim
 \frac{\beta^{k+1}}{2^{9/4}\pi^{3/2}}\,
 \frac{\lambda^n}{n^{3/2}},\label{eq:Tasy}\\
 U(n,n+k)&\sim
 \frac{\varphi^{2k}\rho^{1/2}}{2\pi 5^{1/4}}\,
 \frac{\rho^n}{n},\label{eq:Upasy}\\
 U(n+k,n)&\sim
 \frac{\varphi^{3k}\rho^{1/2}}{2\pi 5^{1/4}}\,
 \frac{\rho^n}{n}.\label{eq:Umasy}
\end{align}
\end{theorem}

\subsection{The discrete Laplace argument, with endpoint control}
We record explicitly the analytic step used to prove the theorem.
Suppose a positive summand, uniformly for $t=j/n$ in each closed
subinterval of $(0,1)$, has the form
\[
 w_{n,j}=n^{-a}A(t)\exp(n\Phi(t))(1+O(n^{-1})),
\]
where $A$ is continuous and positive near a unique maximizer
$\tau\in(0,1)$ of $\Phi$, and $\Phi''(\tau)<0$.
Suppose also that, outside any fixed neighborhood of $\tau$, the
summands have a uniform upper bound
$n^C\exp(n\Phi(t))$ for some fixed $C$. Then
\begin{equation}\label{eq:laplace}
 \sum_{j=0}^n w_{n,j}\sim
 n^{1/2-a}A(\tau)\exp(n\Phi(\tau))
 \sqrt{\frac{2\pi}{-\Phi''(\tau)}}.
\end{equation}

Here is a justification that also covers the applications below.
Continuity and uniqueness of the maximum give an exponential gap
outside any fixed neighborhood of $\tau$; the factor $n^C$ and the
$n+1$ possible summands cannot overcome that gap.
In a sufficiently small fixed neighborhood there is a $c>0$ such that
$\Phi(t)\leq\Phi(\tau)-c(t-\tau)^2$.
Writing $j=n\tau+u\sqrt n$, the central portion is a Riemann sum with
mesh $n^{-1/2}$. Taylor's formula gives, for each bounded $u$ interval,
\[
 n(\Phi(\tau+u/\sqrt n)-\Phi(\tau))
 \longrightarrow\tfrac12\Phi''(\tau)u^2.
\]
The quadratic upper bound supplies an integrable Gaussian majorant.
Truncate first to bounded $u$, pass to the Riemann-sum limit, and then
let the truncation grow. The limiting integral is
$\sqrt{2\pi/(-\Phi''(\tau))}$, proving \eqref{eq:laplace}.

For our factorial summands, the required global polynomial upper bound
follows from
\[
 \log(r!)=r\log r-r+O(\log(r+2))\quad(0\leq r\leq 2n+k),
\]
with the convention $0\log0=0$. There are only finitely many factorial
factors, so the total logarithmic error is $O_k(\log(n+2))$ even at
$j=0,n$. This is a consequence of Stirling's formula; using only the
coarser logarithmic form at the endpoints avoids any division by $t$
or $1-t$ there. On a compact interior interval, the usual square-root
prefactors and relative error $1+O(n^{-1})$ are uniform.

\subsection{Proof of the cubic asymptotic}
The summand on the main diagonal is
\[
 c_{n,j}=\frac{(n+j)!^2}{(n-j)!^2j!^4}.
\]
For $t=j/n\in(0,1)$ define
\[
 \Phi_3(t)=2\bigl((1+t)\log(1+t)
 -(1-t)\log(1-t)-2t\log t\bigr).
\]
Stirling's formula gives, uniformly on compact interior intervals,
\begin{equation}\label{eq:stirling3}
 c_{n,j}=\frac{1+t}{4\pi^2(1-t)t^2}\,
 n^{-2}\exp(n\Phi_3(t))(1+O(n^{-1})).
\end{equation}
The derivative is
\[
 \Phi_3'(t)=2\log\frac{1-t^2}{t^2}.
\]
It decreases from positive to negative and vanishes only at
$\tau_3=1/\sqrt2$. At that point,
\[
 \Phi_3''(\tau_3)=-8\sqrt2,\qquad
 \exp(\Phi_3(\tau_3))=\lambda,
\]
and the prefactor in \eqref{eq:stirling3}, without $n^{-2}$, equals
$\beta/(2\pi^2)$.

For the shifted array $T(n,n+k)$, the ratio of its summand to
$c_{n,j}$ is exactly
\begin{equation}\label{eq:shiftratioT}
 \prod_{\ell=1}^k\frac{n+j+\ell}{n-j+\ell}.
\end{equation}
On a compact interior interval this converges uniformly to
$((1+t)/(1-t))^k$, whose value at $\tau_3$ is $\beta^k$.
For fixed $k$ it is at most polynomially large in $n$, including at
$j=n$, so the endpoint argument preceding \eqref{eq:stirling3} still
applies. Formula \eqref{eq:laplace}, with $a=2$, therefore gives
\[
 T(n,n+k)\sim n^{-3/2}\lambda^n\,
 \frac{\beta^{k+1}}{2\pi^2}
 \sqrt{\frac{2\pi}{8\sqrt2}}
 =\frac{\beta^{k+1}}{2^{9/4}\pi^{3/2}}
 \frac{\lambda^n}{n^{3/2}}.
\]
This proves \eqref{eq:Tasy}.

\subsection{Proof of the quadratic asymptotics}
The main-diagonal summand for $U$ is
\[
 d_{n,j}=\frac{n!(n+j)!}{j!^3(n-j)!^2}.
\]
Define
\[
 \Phi_2(t)=(1+t)\log(1+t)-2(1-t)\log(1-t)-3t\log t.
\]
Then
\begin{equation}\label{eq:stirling2}
 d_{n,j}=(2\pi n)^{-3/2}
 \frac{\sqrt{1+t}}{t^{3/2}(1-t)}
 \exp(n\Phi_2(t))(1+O(n^{-1}))
\end{equation}
uniformly on compact interior intervals. Its saddle equation is
\[
 \Phi_2'(t)=\log\frac{(1+t)(1-t)^2}{t^3}=0.
\]
Since $(1+t)(1-t)^2-t^3=1-t-t^2$, the unique maximum is at
$\tau_2=1/\varphi$. Elementary substitution gives
\[
 \exp(\Phi_2(\tau_2))=\rho,\qquad
 -\Phi_2''(\tau_2)=4\varphi+3=\sqrt5\,\varphi^3.
\]
Using \eqref{eq:laplace} in \eqref{eq:stirling2} yields the constant
\[
 \frac1{2\pi}
 \frac{\sqrt{1+\tau_2}}
 {\tau_2^{3/2}(1-\tau_2)\sqrt{-\Phi_2''(\tau_2)}}
 =\frac{\varphi^{5/2}}{2\pi5^{1/4}}
 =\frac{\rho^{1/2}}{2\pi5^{1/4}}.
\]

The ratio of a summand of $U(n,n+k)$ to $d_{n,j}$ is
\[
 \prod_{\ell=1}^k\frac{n+\ell}{n-j+\ell}
 \longrightarrow (1-t)^{-k},
\]
which equals $\varphi^{2k}$ at the saddle. For $U(n+k,n)$ the ratio is
\[
 \prod_{\ell=1}^k\frac{n+j+\ell}{n-j+\ell}
 \longrightarrow\left(\frac{1+t}{1-t}\right)^k,
\]
with saddle value $\varphi^{3k}$.
Both ratios are polynomially bounded at the endpoints for fixed $k$.
The same Laplace argument therefore proves \eqref{eq:Upasy} and
\eqref{eq:Umasy}, completing the proof of
Theorem~\ref{thm:arrayasy}.

\subsection{Exact leading constants for the remainders}
\begin{theorem}[Sharp shifted-diagonal errors]\label{thm:errorasy}
For each fixed $k\geq0$,
\begin{align}
 \zeta(3)-\Sthree{N}{k}
 &\sim4\pi^3\lambda^{-(2N+k+1)},\label{eq:errasy3}\\
 \zeta(2)-\Stwop{N}{k}
 &\sim(-1)^N4\pi^2\varphi^{-4k}\rho^{-(2N+1)},\label{eq:errasy2p}\\
 \zeta(2)-\Stwom{N}{k}
 &\sim(-1)^{N+k}4\pi^2\varphi^{-6k}\rho^{-(2N+1)}.\label{eq:errasy2m}
\end{align}
For signed quantities, $x_N\sim y_N$ means $x_N/y_N\to1$.
\end{theorem}
\begin{proof}
Denote the positive cubic summand in \eqref{eq:main3} by $d^{(3)}_{n,k}$.
Since its explicit rational prefactor is asymptotic to $6/n^3$,
\eqref{eq:Tasy} gives
\begin{equation}\label{eq:term3asy}
 d^{(3)}_{n,k}\sim96\sqrt2\,\pi^3\lambda^{-k}\lambda^{-2n}.
\end{equation}
For the magnitudes of the superdiagonal and subdiagonal quadratic
summands, \eqref{eq:Upasy} and \eqref{eq:Umasy} similarly give
\begin{align}
 d^{(2,+)}_{n,k}&\sim20\sqrt5\,\pi^2\varphi^{-4k}\rho^{-2n},\label{eq:term2pasy}\\
 d^{(2,-)}_{n,k}&\sim20\sqrt5\,\pi^2\varphi^{-6k}\rho^{-2n}.\label{eq:term2masy}
\end{align}

To justify summing these equivalents, if $d_n=Cq^n(1+o(1))$ with
$0<q<1$, then the relative error is uniformly small for \emph{all}
$n>N$ once $N$ is large. Hence the absolute difference between its tail
and the corresponding geometric tail is at most
$\varepsilon Cq^{N+1}/(1-q)$. This proves
\[
 \sum_{n>N}d_n\sim\frac{Cq^{N+1}}{1-q},\qquad
 \sum_{n>N}(-1)^{n+1}d_n
 \sim(-1)^N\frac{Cq^{N+1}}{1+q}.
\]
For the alternating formula, the error relative to the geometric tail
is at most $\varepsilon(1+q)/(1-q)$, which also tends to zero.
Thus no informal interchange of an asymptotic limit and an infinite
sum is required.

Apply these facts with $q=\lambda^{-2}$ or $q=\rho^{-2}$.
The simplifications
\[
 \lambda^2-1=24\sqrt2\,\lambda,\qquad
 \rho^2+1=5\sqrt5\,\rho
\]
then give exactly \eqref{eq:errasy3}--\eqref{eq:errasy2m}.
\end{proof}

The asymptotic decimal accuracy gained per additional diagonal term is
therefore
\[
 2\log_{10}\lambda=3.0622054827029046\ldots
 \quad\text{for }\zeta(3),
\]
and
\[
 2\log_{10}\rho=2.0898764024997873\ldots
 \quad\text{for }\zeta(2).
\]
Increasing a \emph{fixed} shift improves the leading constant by the
explicit factors shown in Theorem~\ref{thm:errorasy}. It does not change
the per-term exponential rate. These asymptotic formulas describe
performance but are not substituted for the rational certificates in
Theorems~\ref{thm:error3} and \ref{thm:error2}.

\section{Exact verification and reproducible certificates}\label{sec:computation}

\subsection{What was checked}
The accompanying \texttt{verify.py} uses only Python's standard library:
integer binomial coefficients and exact reduced rational numbers.
Its run produced 25,840 successful checks. Among them are independent
binomial expressions for both arrays on $0\leq n,m\leq24$, all four
contiguity identities on $1\leq n,m\leq48$, square closure, termwise
certificate identities including their endpoints, 480 finite identities
for each of the three diagonal families, 300 random monotone paths for
each potential, row and shifted recurrences, and interval consistency.
The pseudorandom seed is fixed at 20260920.

The main verification does not call a numerical zeta function.
A separate optional \texttt{check\_symbolic.py} uses SymPy to simplify
seven rational-function identities to zero. Another optional program,
\texttt{diagnostics.py}, compares array and error asymptotics with
420-digit mpmath evaluations; those numerical comparisons are explicitly
labeled diagnostics, not certificates or proofs.

Finite tests cannot establish the infinitely quantified theorems.
Their purpose is to expose transcription, indexing, sign, or
implementation errors independently of the written derivations.
The identities in Sections~\ref{sec:cubic}--\ref{sec:proof2} are the
proofs for all parameters.

\subsection{How an exact interval certifies decimal digits}
For example, compute $s=F_3(N,N+k)$ exactly and set
\[
 B=\frac1{2(N+k)^2T(N,N+k)^2}.
\]
Then $\zeta(3)\in(s,s+B)$. For $\zeta(2)$, use the smaller rational
bound from Corollary~\ref{cor:diag2}, and put the interval on the correct
side of the partial sum according to its proven error sign.
All interval endpoints are rational numbers whose full numerators and
denominators are saved in \texttt{certificates.json}.

For an interval $[L,R]$ and a positive integer $d$, the exact condition
\begin{equation}\label{eq:decimalcert}
 \lfloor10^dL\rfloor=\lfloor10^dR\rfloor
\end{equation}
certifies the first $d$ decimal digits after the point of every number
in the interval. The program checks \eqref{eq:decimalcert} using integer
division, not floating-point rounding. The following counts are the
actual common-prefix lengths of the generated rational intervals;
they are not inferred from the leading asymptotic formulas.

\begin{center}
\begin{tabular}{@{}rccc@{}}
\toprule
Shift $k$ & $\zeta(3)$ & $\zeta(2)$ superdiagonal & $\zeta(2)$ subdiagonal\\
\midrule
$0$  & $180$ & $124$ & $124$\\
$1$  & $182$ & $125$ & $125$\\
$3$  & $185$ & $127$ & $128$\\
$10$ & $194$ & $131$ & $136$\\
\bottomrule
\end{tabular}
\end{center}
\noindent All entries use $N=60$ summands after the boundary harmonic sum.
An interval may have a very small width but cross a decimal boundary;
this is why the program checks common digits rather than merely
reporting $-\log_{10}$ of the width.

The archive also contains \texttt{convergence\_table.csv},
\texttt{array\_sample.csv}, complete verification reports, a source audit,
and build instructions. The generated values are reproducible from
\texttt{verify.py}; they are not needed to trust any numerical oracle.

\section{What the proof establishes, and what it does not}\label{sec:scope}

The main conclusion is exact and complete: the entire shifted-diagonal
family for $\zeta(3)$ in A143007 and both shifted-diagonal families for
$\zeta(2)$ in A108625 hold for every nonnegative integer shift.
The key reduction is also explicit. A finite hypergeometric sum gives
a telescoping certificate; the certificate gives a four-corner identity;
that identity makes a rational lattice sum path-independent; and the
boundary values determine its limit. Summing along a diagonal then
produces the polynomial numerator appearing in the conjecture.

Several stronger consequences are established within the same argument:
exact horizontal and vertical remainders, convergence along arbitrary
unbounded monotone lattice paths, signed rational bounds, shifted cubic
recurrences, and sharp asymptotic error constants. All limiting steps
are supplied after, rather than assumed from, the local algebra.

The following distinctions matter. This work does not settle the
irreducibility conjectures in the classical \Apery{} entries or the
other congruence conjectures listed in the two array entries. It does
not claim a new proof of irrationality from fast convergence alone.
It is not a proof-assistant formalization, and the accompanying finite
checks do not replace mathematical proof. Finally, because related
conservative fields and their conjugate polynomials already occur in
~\cite{david}, no historical novelty or priority claim is made for the
framework or for consequences that might already be implicit in that
literature. The specific accomplishment here is an explicit,
self-contained derivation of the three formulas that the consulted
OEIS entries still label conjectural.

\appendix
\section{A compact route through the proof}
A reader interested only in the conjectural identities can omit the
recurrence and asymptotic sections. The logically sufficient chain is
\[
 \eqref{eq:deltaW},\ \eqref{eq:deltaV}
 \quad\Longrightarrow\quad
 \eqref{eq:contig3},\ \eqref{eq:contig2}
 \quad\Longrightarrow\quad
 \eqref{eq:closed3},\ \eqref{eq:closed2}.
\]
Lemma~\ref{lem:closed} then constructs $F_3,F_2$; the elementary boundary
limits prove their convergence; and the diagonal increments
\eqref{eq:step3}, \eqref{eq:step2} telescope to the three theorems.
The sum certificates vanish at both endpoints, all corner denominators
are positive, and equal coordinates are permitted throughout.

\section{Rebuilding the supplementary artifacts}
From the directory containing the archive's files, run
\begin{verbatim}
python3 verify.py
pdflatex -interaction=nonstopmode -halt-on-error article.tex
pdflatex -interaction=nonstopmode -halt-on-error article.tex
\end{verbatim}
The main program needs Python 3.10 or later and no third-party Python
packages. The optional checks are
\begin{verbatim}
python3 check_symbolic.py   # requires SymPy
python3 diagnostics.py      # requires mpmath
\end{verbatim}
No script accesses the network. The supplied PDF was compiled from
this source and rendered for visual inspection. Bibliographic links
identify the original sources; the archive does not redistribute
complete third-party articles or OEIS pages.

\begin{thebibliography}{9}
\small
\setlength{\itemsep}{3pt}
\bibitem{oeisT}
OEIS Foundation Inc.,
\emph{A143007: Square array related to the crystal ball sequences
for the product lattices $A_n\times A_n$},
\url{https://oeis.org/A143007}.
Consulted September 20, 2026; formula field, conjectural shifted diagonals.

\bibitem{oeisU}
OEIS Foundation Inc.,
\emph{A108625: Crystal ball sequence for the root lattices $A_n$},
\url{https://oeis.org/A108625}.
Consulted September 20, 2026; formula field, conjectural superdiagonals
and subdiagonals.

\bibitem{oeis258}
OEIS Foundation Inc.,
\emph{A005258: Ap\'ery numbers},
\url{https://oeis.org/A005258}.
Consulted September 20, 2026.

\bibitem{oeis259}
OEIS Foundation Inc.,
\emph{A005259: Apery (Ap\'ery) numbers},
\url{https://oeis.org/A005259}.
Consulted September 20, 2026.

\bibitem{apery}
R. Ap\'ery,
\emph{Irrationalit\'e de $\zeta(2)$ et $\zeta(3)$},
Ast\'erisque \textbf{61} (1979), 11--13.
Bibliographic record linked from OEIS A005258 and A005259.

\bibitem{david}
O. David,
\emph{The conservative matrix field},
arXiv:2303.09318v3 (2023), revised December 4, 2023.
\url{https://arxiv.org/abs/2303.09318}.
The HTML version consulted is
\url{https://arxiv.org/html/2303.09318v3}.

\bibitem{straub}
A. Straub,
\emph{Multivariate Ap\'ery numbers and supercongruences of rational functions},
Algebra \& Number Theory \textbf{8} (2014), no.~8, 1985--2008.
\href{https://doi.org/10.2140/ant.2014.8.1985}{doi:10.2140/ant.2014.8.1985}.
Preprint: \url{https://arxiv.org/abs/1401.0854}.
\end{thebibliography}
\end{document}
